\documentclass{article}

% Language setting
% Replace `english' with e.g. `spanish' to change the document language
\usepackage[english]{babel}
\usepackage{placeins}  % 导言区


% Set page size and margins
% Replace `letterpaper' with `a4paper' for UK/EU standard size
\usepackage[letterpaper,top=2cm,bottom=2cm,left=3cm,right=3cm,marginparwidth=1.75cm]{geometry}
\usepackage[affil-it]{authblk} % 加载 authblk 宏包

% Useful packages
\usepackage{amsmath}
\usepackage{graphicx}
\usepackage[colorlinks=true, allcolors=blue]{hyperref}

% 提及无监督的概念，
\title{Representation Learning on Comprehensive Neuron Skeletons for Micro-Connectome Analysis}

\author{
}

\date{
    $^1$ Institution A \\
    $^2$ Institution B \\
    $^3$ Institution C \\
}


\begin{document}
\maketitle
\begin{abstract}
Serial electron microscopy provides micro-connectome datasets containing massive neuronal morphologies and connectivity maps, posing analytical challenges.  
Automated methods such as NBLAST, which employs curve-distance algorithms, and GraphDINO, which applies contrastive learning techniques, accelerate skeleton-based neuronal analysis but cannot simultaneously achieve%multimodal integration.  
%Here, by assigning features such as diameters, synaptic targets or any local node embedding to fine skeleton nodes, comprehensive skeletons can maximize the utility of micro-connectomic data at relatively low computational cost.  
However, conventional graph neural networks remain insufficient for contrastive learning on fine skeletons with large numbers of nodes, a limitation we address using subgraph decomposition combined with a tailored hierarchical graph neural network architecture (HGNNA).  
We apply contrastive learning to the FlyWire dataset using simplified skeletons with 200 nodes and comprehensive skeletons with 3500 nodes, and testing on Kenyon Cell (KC) classification shows that the richer neuronal information in comprehensive skeletons can be captured by graph neural network to improve prediction accuracy. Subsequently, in the absence of cell type labels, the morphology embeddings of KC neurons and their axonal segments independently show that Hemilineage and cell types are significantly separated under t-SNE dimensionality reduction. Thereafter, the testing results on the FlyWire optic lobe and MICrONS datasets are less straightforward, and the article discusses the relationship between contrastive learning, morphology, and absolute coordinates in terms of input, suggesting that appropriate input features and data augmentation strategies should be adopted across different species or brain regions. Finally, we reassign the learned morphological representations of neurons in FlyWire as embedding vectors to their associated synaptic nodes. The article thus discusses the compatibility of the method with embedding vectors and the potential extension to connectivity analysis.
\end{abstract}

\section{Introduction}
% 【电镜数据——发展需求】 % *背景和重要性
Understanding how the brain operates requires a detailed characterization of its underlying wiring architecture. Connectomics, which aims to map neural circuits at synaptic resolution, provides a fundamental structural blueprint for revealing these principles. 
In recent years, advances in serial-section electron microscopy (ssEM) have enabled the reconstruction of large-scale, nanometer-resolution neural morphologies and circuits, yielding unprecedented datasets for connectomic research (the Drosophila hemibrain, FlyWire, and the MICrONS project). These connectome offer detailed neuronal reconstructions and synaptic connectivity for up to ~140,000 neurons (e.g., FlyWire). 
Before a neuronal system can be directly modeled from connectome and simulated in silico, neurons are systematically categorized into cell classes and organized into meaningful clusters to generate an interpretable network map. However, deriving biologically meaningful categories from dense connectomic remains a major challenge. The complexity of synaptic networks, the high morphological diversity of neurons, and the scale of modern ssEM datasets make manual annotation both time-consuming and heavily dependent on expert knowledge.

% 【电镜数据——数据形式、现有研究——形态、连接】*现有研究的局限性 *研究趋势 *在新趋势下的关键问题
Historically, neuronal analysis relied on manually defined statistical features. Although supervised deep learning and unsupervised clustering improved accuracy and efficiency, they remain constrained by a reliance on hand-crafted descriptors and expert validation, which limits scalability and introduces bias. 
In practice, neuronal types in previously unexplored neural systems are typically inferred by comparing and summarizing morphological and connectivity patterns across large populations of neurons. Even when it is known which features are informative for neuronal classification, extracting and validating these features requires extensive and careful analysis by domain experts. Hence, our work is motivated by a fundamental premise: if our understanding of a target neural system is extremely limited—analogous to the perspective of a non-expert numerical analyst—can meaningful organizational principles still be discovered directly from the data, enabling the identification of similar neuronal types and the discrimination of distinct ones?
To mitigate these issues, unsupervised representation learning has been introduced into this area, utilizing contrastive learning and Graph Neural Networks (GNN) to analyze connectome data from a uniformly modified initial form. These approaches typically treat morphology as skeletons or connectivity as connectivity graph separately. Crucially, however, insights from circuit discovery indicate that neither feature is sufficient in isolation; rather, the specific connectivity along neuron neurites is essential for precise neuronal identification.


% 【神经元的精细建模】 *我们的解决方案，精细骨架及挑战* 不提融合图谱的概念，而是从单细胞视角加入连接性这一概念。To enrich the neuritic skeleton是为了实现对应的表示学习
To enrich the neuritic skeleton with both structural and connection-specific attributes, a proper initial form of finely sampled skeletons with extra node feature is proposed. A high-resolution resampling of the neuritic tree that increases node density enables precise mapping of individual synapses onto the underlying morphology. However, this dense representation substantially increases the computational burden, making large-scale graph learning prohibitively expensive. % 【子图分解和分级图神经网络框架】 *我们的解决方案子图分解和HGNNA
To overcome this limitation, we introduce a skeleton decomposition strategy in conjunction with a hierarchical graph neural network architecture (HGNNA). This approach reduces the computational cost by more than an order of magnitude, enabling representation learning for hundreds of thousands of neurons with fine structure and connectivity under practical resource constraints.

% 【空间对称性在不同系统中存在差异】 *问题：坐标
During our investigation, we found that neuronal coordinate systems exhibit distinct symmetry properties across biological systems. For instance, cortical neurons of the same type often display translational symmetry, whereas typical neurons in Drosophila instead follow mirror symmetry. Within the fly optic lobe, yet another pattern emerges: neurons exhibit symmetry consistent with movement along a curved surface. These heterogeneous coordinate symmetries impose inherent constraints on methods that explicitly rely on spatial positions. In contrast, when spatial coordinates are excluded, the topological organization of neuronal skeletons may remain largely consistent across neurons of similar types, irrespective of the symmetry of their underlying coordinate systems.

% 【形态、拓扑、连接】 *具体研究的路线
Motivated by this observation, we construct skeletons with coordinates as morphological representations and skeletons without coordinates as topological representations. Connectivity information is further incorporated by enriching each skeleton node with synaptic annotations and target embeddings. This joint formulation enables a unified analysis of both neuronal morphology and connectivity at the level of individual cells.
% 【FlyWire和MICrONS】 *具体研究应用的点
The validity and generality of the proposed representations are assessed using the FlyWire and MICrONS datasets. For the FlyWire dataset, we adopt the classification of Kenyon cell (KC) neuron populations as the main benchmark to assess the representation capability.
% 为了将大型骨架的图神经网络学习我们从拓扑和连接两个角度【大型骨架和模型】【拓扑结构】【拓扑连接性】【应用】 

% **精细采样骨架的无监督学习为分类任务捕获更多信息**
\section{Unsupervised Learning on Finely Sampled Skeletons Captures Richer Information for Classification Tasks} % 更多的描述放在method中
% 【骨架节点数统计】
To address the question of how many nodes are sufficient for the description of a neuron skeleton, we selected a publicly available dataset: FlyWire (adult Drosophila) for statistical analysis. FlyWire provide finely sampled skeletons and detailed synaptic information. For each synapse, we assigned it to the nearest node of the corresponding skeleton and labeled the nodes as either 'pre-synapse,' 'post-synapse,' or 'fork node,' and 'end node' where appropriate (a demostration is shown in {Fig.\ref{fig:morph_classify}a}). A node is considered a "meaningful node" if it is labeled as one of these four types. We then counted the number of meaningful nodes for each skeleton in the dataset. The distribution of meaningful node counts across all skeletons is shown in {Fig.\ref{fig:morph_classify}b}. Most skeletons containing fewer than 1,000 meaningful nodes, while a small fraction possesses several thousand nodes. Based on the cumulative percentage calculation ({Fig.\ref{fig:morph_classify}c}) for the distribution of meaningful nodes, we set the threshold at 3,500 nodes in our application, which covers 96.37\% of cells in FlyWire. Skeletons with more than 3500 nodes will have nodes randomly dropped to maintain a skeleton with 3500 nodes.

% Figure 1. 骨架节点数统计 # MICrONS可以放在extended
\begin{figure}[htbp]
\centering
\includegraphics[width=1\linewidth]{figures/Fig1_0202.pdf}
\caption{(a) Example of a neuron with end, fork, pre- or post-synapse label. (b) Distribution of meaningful node counts across all neuron skeletons in FlyWire. (c) Cumulative percentage of neuron skeletons as a function of meaningful node counts. (d) t-SNE visualization of morphology representations colored by super classes for all neurons in FlyWire. (e) Pie graph of the target neurons groups of four tasks. (f) Schematic of cell-type classification after contrastive learning of morphology representations. (g) Classification accuracy comparison between coarse and finely sampled skeletons across four tasks: Central Brain (CB), Optic Lobe (OL), Kenyon Cells (KC), and Transmedullary neurons (Tm). (h) Example morphologies of KC and Tm neurons.}
\label{fig:morph_classify}
\end{figure} 

% 3500 进行 单层8 节点子图分解，邻接矩阵大小将会变为原来的1.8%。
% 【测试的方法和对象】
We implemented a skeleton decomposition combined with a hierarchical graph neural network (HGNNA) enabling contrastive learning on finely sampled neuron skeletons. Through contrastive learning process, the model learns to extract intrinsic morphology differences directly from the data. After optimization, each input is mapped to a latent embedding that encodes rich information about the underlying neuronal morphology, so called morphology representation. 
To verify whether the enhanced expressiveness of comprehensive skeletons can be effectively utilized by our designed HGNNA, we conducted cell-type classification tasks on the FlyWire dataset. Following the morphology representation learning stage, the projection head was replaced with a classification head for supervised evaluation ({Fig.\ref{fig:morph_classify}f}). The test is performed on the classification of types in central brain (CB), optic lobe (OL), Kenyon cells (KC), and transmedullary neurons (Tm) types. Among these, CB and OL constitute a large proportion (38.9 percent) of all neurons in FlyWire, representing central and peripheral neural populations, respectively; we performed cell-type classification tests on these two populations. KC and Tm are neuronal classes in CB and OL, respectively, and can be further subdivided into more detailed types, with neuronal morphologies and distributions shown in {Fig.\ref{fig:morph_classify}h} and {Fig.\ref{fig:morph_sensitivity}g}. 


% 【测试的结果】不是都采用缩写是为了提高可读性。KC和Tm采用了缩写。
The classification accuracy results are summarized in {Fig.~\ref{fig:morph_classify}g}. Across all four classification tasks, finely sampled skeletons consistently outperform coarse skeletons. Specifically, the accuracy improvements are: Central brain classification improves from 75.38\% to 90.81\% (15.43\% gain), optic lobe classification from 58.72\% to 89.71\% (30.99\% gain), KC classification from 81.74\% to 90.00\% (8.26\% gain), and Tm classification from 70.03\% to 100.00\% (29.97\% gain). 
Notably, while the visual system neurons (in the tasks of OL and Tm) demonstrate notably larger accuracy gains compared to central brain populations, this pronounced improvement in visual system classification suggests that neuronal types in the optic lobe are more tightly coupled with their morphological features. The exceptional performance on Tm neurons (100\% accuracy with finely sampled skeletons) indicates that morphology representation alone can completely distinguish Tm types, contrary to previous studies suggesting that some Tm types require connectivity information for proper classification. This demonstrates that the detailed morphological information captured by fine-grained skeleton representations is particularly effective for distinguishing visually-specialized neuronal types.


% **形态对非功能性特征、空间位置和对称性敏感**
\section{Morphology Representation is Sensitive to Nonfunctional Features, Spatial Position, and Symmetry} % 主要
% 转折
% 用KC说明影响因素：发育谱系、对称性，神经元之间的空间竞争。
% 说明形态表示的局限性。
% 现在的图内容还可以，这个后续再优化，先放下不动。
% 形态-Tm的结果放在这里一部分。从这一点出发，我们将FlyWire神经元骨架经过GNN后输出的latent表示进行t-SNE降维。为了观察不同类别神经元的latent分布情况，将不同类型神经元用不同颜色表示（图1d 所有FlyWire神经元不同super class的分布情况，图2a 所有KC神经元不同类型的分布情况，图2h Tm不同类型的分布情况）。

It is worth noting that contrastive learning is an unsupervised method for deriving features from large datasets. Thus, under the assumption of having no cell type labels, one may ask whether the latent representations can be exploited to analyze the commonalities and differences among neurons. Starting from this perspective, we performed t-SNE dimensionality reduction on the latent representations output by the GNN for FlyWire neuron skeletons. To visualize the latent distribution of different neuronal categories, neurons of different types were represented using different colors ({Fig.~\ref{fig:morph_classify}d} shows the distribution of different super classes across all FlyWire neurons, {Fig.~\ref{fig:morph_sensitivity}a} shows the distribution of different types among all KC neurons, and {Fig.~\ref{fig:morph_sensitivity}h} shows the distribution of different Tm types).

% Figure 2. 
\begin{figure}[htbp]
\centering
\includegraphics[width=1\linewidth]{figures/Fig2_0202.pdf}
\caption{(a) and (b) t-SNE visualization of morphology representations of Kenyon Cells (KCs). (c) Morphological differences in dendritic segments among KC hemilineages. (d) t-SNE visualization of morphology representations of KCs after removing dendritic segments. (e) Example of a misclassified KC neuron. (f) Pie chart showing the correspondence between clustering types and original label types. (g) Morphological differences in Tm neurons along the y-coordinate. (h) t-SNE visualization of morphology representations of Tm neurons. (i) Correlation between the y-coordinate of the soma and the latent representation of Tm neurons.}
\label{fig:morph_sensitivity}
\end{figure}

% 错误类型标签存在于calssificaiton.csv中，网页上直接查询是正常的。
The distribution of KC types exhibits multiple clusters ({Fig.~\ref{fig:morph_sensitivity}b}). Multiple clusters correspond to KC cell types, different hemilineages and left-right brain hemispheres. KC neurons display pronounced morphological differences in their dendritic segments due to distinct developmental lineages ({Fig.~\ref{fig:morph_sensitivity}c}), and these differences are independent of KC classification types (ref.X). By removing dendritic segments from the network input without retraining, we can eliminate the effects of developmental lineage ({Fig.~\ref{fig:morph_sensitivity}d}), revealing only left-right hemispheric and KC cell type differences. Clustering the morphological representations without dendritic segments yields the primary classification clusters within KC types; under this five-type clustering scheme, consistency reaches X\% (the correspondence between clustering types and original label types is shown in the pie chart in {Fig.~\ref{fig:morph_sensitivity}f}). Furthermore, by examining neurons of other types that fall within single-type clusters in {Fig.~\ref{fig:morph_sensitivity}d}, we identified instances of mislabeling. {Fig.~\ref{fig:morph_sensitivity}e} shows a KCa'b' neuron that was incorrectly classified as KCg-m in the type labels.

% !补充Tms中的latent与y坐标关联性
The t-SNE visualization of Tm neurons reveals a more dispersed distribution, with Silhouette score of 0.008. This fragmentation is partially attributable to the rotational and translational symmetries exhibited by optic lobe neurons, as well as the presence of non-functional segments arising from developmental processes. By computing the correlation between the y-coordinate of the soma and the latent representation of the cell, a significant association emerges, as visualized in {Fig.~\ref{fig:morph_sensitivity}i}. 
These findings indicate that non-functional segments arising from development and the spatial symmetry of neurons significantly impact morphological representations, posing challenges for unsupervised clustering.

% **拓扑表示在聚类方面的优势**
\section{Topology Representation Excels in Clustering} 

A portion of the sensitivity of morphological representations arises from their dependence on spatial coordinates. To disentangle this effect, we removed coordinate inputs while retaining the skeletal topological structure, which we refer to as the topological representation (Fig.\ref{fig:topology}a). Compared with the morphological representation, the topological representation exhibits consistently improved intra-class compactness and inter-class separability across multiple evaluation settings, including all-neuron super classes, optic-lobe classes, Tm classes, and KC classes (Fig.\ref{fig:topology}b–e, g–i). 

% 四个任务具体的数值
Specifically, for KC neurons, the average inter-class distance increases from 5.639 (morphology) to 8.382 (topology), representing a 1.49-fold improvement in class separation, while the average intra-class distance decreases from 4.900 to 3.721, indicating 24.1\% tighter within-class clustering. The Silhouette score improves from -0.048 (morphology) to 0.090 (topology). For Tm neurons, the inter-class distance increases from 12.077 to 12.692 (1.05-fold), while intra-class distance decreases from 9.445 to 6.696 (29.1\% reduction), with Silhouette scores of 0.008 and 0.210 respectively. Similar improvements are observed for neurons in CB (Silhouette scores: -0.019 vs 0.066) and OL (Silhouette scores: -0.149 vs 0.059). Together, these quantitative metrics demonstrate that the topological representation (Fig. \ref{fig:topology}e) provides clear advantages over the morphological representation (Fig.Fig.~\ref{fig:morph_sensitivity}d) in unsupervised clustering.

% 小比例的OL Tm CB重新测试小训练集分类
% These advantages are further reflected in few-label classification performance. We conducted classification experiments using only a small fraction of labeled data, selecting limited proportions of labels for CB, OL, KC, and Tm classification tasks. As shown in Fig. \ref{fig:topology}f, the topological representation consistently outperforms the morphological representation across all four tasks, achieving accuracy improvements of X1%, X2%, X3%, and X4%, respectively.
These advantages are further reflected in few-label classification performance. We conducted classification experiments using only a small fraction of labeled data, selecting limited proportions of labels for CB, OL, KC, and Tm classification tasks. As shown in Fig. \ref{fig:topology}f, in the five-class Kenyon cell classification task, reducing the labeled data ratio from 80\% to 10\% leads to a pronounced performance degradation in morphology-based representations (from 98\% to 85\%). By contrast, the accuracy of topology-based classification remains nearly unchanged and exceeds that of morphology-based classification at 10\% labeling, highlighting its improved robustness.

% Figure 3. 形态聚类尝试与局限性,KC Tm
\begin{figure}[htbp]
\centering
\includegraphics[width=1\linewidth]{figures/Fig3_0202.pdf}
\caption{(a) Schematic of topology representation without coordinate inputs. (b–e) t-SNE visualizations of topological representations colored by super classes for all neurons, classes within the optic lobe, Tm classes, and KC classes, respectively. (f) Few-label classification accuracy comparison between topology and morphology representations across four tasks. (g–i) Quantitative metrics of intra-class and inter-class distances for topology representations in all neurons, optic lobe neurons, and Tm neurons, respectively.}
\label{fig:topology}
\end{figure} 




\FloatBarrier

\section{Method}


\subsection{Method: Graph decomposition and hierarchical graph neural network}
% Graph
% Graph decomposition history
% My graph decompositon
% 大型图处理方法有哪些
A graph can be expressed as nodes and edges, and in graph neural network operations, it is commonly represented by feature and adjacency matrices. The complete skeleton of a neuron can contain thousands of nodes to describe its intricate structure, resulting in an adjacency matrix that is particularly large and challenging to process. % Therefore, several methods have been developed, including Sampling Methods (GraphSAGE, PinSAGE), which randomly sample a subset of neighboring nodes to reduce computational complexity; Graph Partitioning (Cluster-GCN), which divides large graphs into multiple subgraphs with nodes feature clustering and executes message passing only within these subgraphs during training; and Graph Pooling (DIFFPOOL), which uses a hierarchical approach that is challenging to train, requiring each level to be trained individually before clustering and inputting to the next level. These methods are primarily applied to complex graphs, such as those formed by social networks, while the skeleton of a single neuron is typically represented as a tree (an acyclic undirected graph) in graph theory, with branches extending from the soma. Therefore, in our work, the method of graph decomposition along with hierarchical graph neural network has been proposed for the efficient processing of neuron skeletons. Our plan has the natural ability of multi-scale feature perception to contain the localized synaptic distribution and global morphology.

% 编辑1
A soma-initiated breadth-first search is performed to partition a tree-like graph into subgraphs, each containing a fixed number of \(n\) nodes. Once a subgraph reaches \(n\) nodes or has no unvisited neighboring nodes, it is finalized. A new unvisited node is then randomly selected as the seed to initiate the next subgraph. The process is visualized in \textbf{Fig}. The decomposition process can be applied multiple times to reduce the adjacency matrix. However, excessive decomposition may lead to fragmentation and increased computational overhead during the decomposition process. In our work, skeletons are decomposed twice to achieve an acceptable memory cost. The adjacency matrix of each subgraph is preserved, and inter-subgraph adjacency matrices are constructed to represent connectivity between subgraphs. A connection between two subgraphs is defined as the existence of at least one edge linking a node in one subgraph to a node in the other.

To facilitate the processing of decomposed subgraphs, a hierarchical graph neural network (HGNN) is designed to extract information from subgraphs without requiring multiple training iterations. Subgraphs with the same number of nodes are batched and processed by a graph neural network (GNN) module, which outputs latent representations for each subgraph. These latent representations along with the inter-subgraph adjacency matrices can be understand as compressed graph of the original graph, and are then processed by another module with the same design to generate a latent representation for a larger subgraph (at a middle layer) or the entire graph (at the final layer). In different hierarchical layers, GNN modules of varying depths are employed to accommodate the differing complexities of information at each scale (normally deeper layers in larger scale). Each GNN module is standardized to simplify the extension of hierarchical layers. Inter-layer adjacency matrices are incorporated so the risk of exploding gradients [\textbf{ref}] can be avoided. As a result, this hierarchical structure enables the GNN to reach significantly greater depths compared to conventional GNNs, thereby leveraging its design to achieve both efficient and scalable learning.

% Memory cost 效率提升定量分析
The critical conflict between an increasing number of skeleton nodes and higher memory cost is significantly alleviated with the application of our method. To quantitatively analyze the reduction in memory cost, we consider a simple case of one-time decomposition where each subgraph contains 
 \(n\) nodes. The memory cost of all adjacency matrices is reduced from \(N^2\) to
\begin{equation}
\label{eq:mcost1}
\frac{N}{n} n^{2} +\left( \frac{N}{n} \right) ^{2},
\end{equation}
where the first term represents the adjacency matrices of multiple subgraphs, and the second term represents the inter-subgraph adjacency matrices. Expression \eqref{eq:mcost1} can be simplified to 
\begin{equation}
\label{eq:mcost1_smp}
N n +\frac{N^{2}}{n^{2}},
\end{equation}
and achieve a minimize value at \(n=\sqrt[3]{2N} \). For example, when a graph has \(N=3500\) nodes, the memory cost is reduced by approximately 122 times when the number of subgraph nodes is set to \(n=19\). However, estimating the most efficient design for the computational model is more complex, as the number of GNN layers may varies across hierarchical levels, and the total number of hierarchical layers may differ. More importantly, instead of prioritizing efficiency alone, subgraphs and compressed graphs must contain sufficient nodes to accurately represent both localized morphology and global morphology.


\subsection{Data Augmentation}

\subsection{GCN}

\subsection{DINO}


Note that your figure will automatically be placed in the most appropriate place for it, given the surrounding text and taking into account other figures or tables that may be close by. You can find out more about adding images to your documents in this help article on \href{https://www.overleaf.com/learn/how-to/Including_images_on_Overleaf}{including images on Overleaf}.


\subsection{How to add Citations and a References List}

You can simply upload a \verb|.bib| file containing your BibTeX entries, created with a tool such as JabRef. You can then cite entries from it, like this: \cite{greenwade93}. Just remember to specify a bibliography style, as well as the filename of the \verb|.bib|. You can find a \href{https://www.overleaf.com/help/97-how-to-include-a-bibliography-using-bibtex}{video tutorial here} to learn more about BibTeX.

If you have an \href{https://www.overleaf.com/user/subscription/plans}{upgraded account}, you can also import your Mendeley or Zotero library directly as a \verb|.bib| file, via the upload menu in the file-tree.

\subsection{Good luck!}

We hope you find Overleaf useful, and do take a look at our \href{https://www.overleaf.com/learn}{help library} for more tutorials and user guides! Please also let us know if you have any feedback using the Contact Us link at the bottom of the Overleaf menu --- or use the contact form at \url{https://www.overleaf.com/contact}.

\bibliographystyle{alpha}
\bibliography{sample}

\end{document}